\scanpage{119}

d'où

\[
\Phi_{n'}(T)^{\varphi(p^r)}(T^n-1)\equiv\deleted{\unsure{\prod}}\,(T^{n'}-1)^{\sum_{s=0}^r\varphi(p^s)}\Phi_n(T)
\]

et comme $\sum_{s=0}^r\varphi(p^s)=1+(p-1)+(p^2-p)+\ldots+(p^r-p^{\unsure{r-1}})$ $=p^r$, on trouve

\[
\Phi_{n'}(T)^{\varphi(p^r)}(T^n-1)\equiv(T^{n'}-1)^{p^r}\Phi_n(T)
\]

Or

\[
T^n-1=(T^{n'})^{p^r}-(1)^{p^r}\equiv(T^{n'}-1)^{p^r}
\]

\deleted{\unsure{comme}} Donc divisant par $(T^{n'}-1)^{p^r}$, on trouve

\[
\Phi_n(T)^{\varphi(p^r)}=\Phi_n(T)
\]

\hfill cqfd.

\uline{Corollaire} Soient $k$ un corps de car. $p>0$, $n$ un entier \deleted{$\geq3,$} $\geq1$, $\alpha\in k$ une racine du polynôme semi-cyclotomique $\Psi_n$, i.e.

\[
\Psi_n(\alpha)=0
\]

Soit \deleted{$r\geq0$ tel que} \inserted{$p^r$ ($r\geq0$) la} plus grande puissance de $p$ qui divise $n$, donc

\[
n=p^rn',\qquad(p,n')=1.
\]

Alors le polygône régulier d'invariant $\alpha$ est d'ordre égal à $n'$ si $n'\geq3$, i.e.

% archive sheet number: 118
